\bookchapter{Κεφάλαιο 5: Κάντε το κινητό να σας κατανοεί — Βασικές ρυθ$\mu$ίσεις}

\chapterintro{Δαπανήστε λίγα λεπτά σε ρυθ$\mu$ίσεις, η ε$\mu$πειρία βελτιώνεται τεράστια. Αυτές οι ρυθ$\mu$ίσεις είναι σαν να «κουρδίζετε» το κινητό, για να ταιριάζει στις συνήθειές σας.}

\booksection{I. Ήχος και δόνηση}

Θέστε ένα αγαπη$\mu$ένο ringtone και notification. Χρήση ευέλικτασίγαση λειτουργία (συνάντηση),$\mu$ην ενοχλείτε λειτουργία (ύπνος), ο κόσ$\mu$ος α$\mu$έσως ησυχάζει.

\begin{itemize}
  \item \textbf{Ringtone}: επιλέξτε ένα που $\mu$πορείτε να ακούσετε και δεν είναι άβολο.
  \item \textbf{Λειτουργία δόνησης}: στη σίγαση ειδοποιεί $\mu$ε δόνηση.
  \item \textbf{Λειτουργία $\mu$η ενόχλησης}: $\mu$πορείτε να ορίσετε εξαιρέσεις, όπως να επιτρέπονται κλήσεις από επαφές.
  \item \textbf{Διαχωρισ$\mu$ός έντασης}: η ένταση πολυ$\mu$έσων, η ένταση κλήσης και η ένταση ξυπνητηριού $\mu$πορούν να ρυθ$\mu$ιστούν ξεχωριστά.
\end{itemize}

\booksection{II. Οθόνη και φωτεινότητα}

\begin{itemize}
  \item \textbf{Αυτό$\mu$ατη φωτεινότητα}: αφήστε το κινητό να ρυθ$\mu$ίζεται αυτό$\mu$ατα βάσει του φωτισ$\mu$ού περιβάλλοντος, άνετο.
  \item \textbf{Λειτουργία προστασίας $\mu$ατιών/νύχτας}: φιλτράρει το $\mu$πλε φως, πιο άνετο το βράδυ.
  \item \textbf{Σκούρα λειτουργία}: εξοικονό$\mu$ηση ενέργειας (ιδιαίτερα αποτελεσ$\mu$ατική σε οθόνες OLED), και δεν ενοχλεί τα $\mu$άτια τη νύχτα. Πολλές εφαρ$\mu$ογές υποστηρίζουν επίσης σκούρα λειτουργία.
  \item \textbf{Ρυθ$\mu$ός ανανέωσης οθόνης}: ορισ$\mu$ένα Android κινητά $\mu$πορούν να ρυθ$\mu$ιστούν, όσο υψηλότερη (π.χ. 120Hz) τόσο πιο ο$\mu$αλή η ολίσθηση.
  \item \textbf{Αυτό$\mu$ατο κλείδω$\mu$α οθόνης}: ρυθ$\mu$ίστε 15 ή 30 δευτερόλεπτα, ασφαλές και εξοικονό$\mu$ηση ενέργειας.
  \item \textbf{Μέγεθος γρα$\mu$$\mu$ατοσειράς}: Σας φαίνονται $\mu$ικρά τα γρά$\mu$$\mu$ατα; Μεγαλώστε τα.
\end{itemize}

\booksection{III. Wallpaper, θέ$\mu$α και επιφάνεια εργασίας}

Αλλάξτε ό$\mu$ορφη wallpaper (στατική ή dynamic), η διάθεση αλλάζει. Τα Android επωνυ$\mu$ίας κινητά έχουν συχνά κατάστη$\mu$α θε$\mu$άτων, αλλάζουν ολόκληρο σύνολο εικονίδια και στυλ.

\textbf{Widgets επιφάνειας εργασίας}: Καιρός, η$\mu$ερολόγιο, χειριστήρια $\mu$ουσικής στην αρχική οθόνη — $\mu$όνο $\mu$ε ένα πάτη$\mu$α. Υποστηρίζεται τόσο σε iOS όσο και σε Android, προσθέστε $\mu$ε παρατετα$\mu$ένο πάτη$\mu$α σε κενό ση$\mu$είο.

\booksection{IV. Ρυθ$\mu$ίσεις αποδοτικότητας: κάντε το κινητό «εξωτερικό εγκέφαλο»}

Αυτή η ενότηταείναι «$\mu$όλις ρυθ$\mu$ιστεί δεν γυρίζετε» λειτουργίες, καθε$\mu$ία$\mu$πορεί να αυξήσει ση$\mu$αντικά την αποδοτικότητά σας.

\booksubsection{Λειτουργία εστίασης/$\mu$ην ενοχλείτε: ανακτήστε την προσοχή που πήραν οι ειδοποιήσεις}

\begin{itemize}
  \item \textbf{Βασική σκέψη}: δεν είναι "κοιτάτε λιγότερο το κινητό ", αλλά "συγκεντρωθείτε όταν πρέπει, επικοινωνήστε όταν χρειάζεται".
  \item \textbf{Συ$\mu$βουλές χρήσης}:\newline · Λειτουργία εργασίας: επιτρέπονται $\mu$όνο επαφές $\mu$ε "αστέρι" και απαραίτητες εφαρ$\mu$ογές·\newline · Λειτουργία ύπνου: ενεργοποιείται αυτό$\mu$ατα, επιτρέπει $\mu$όνο τα ξυπνητήρια.\newline · Γυ$\mu$ναστική / Ανάγνωση / Συνάντηση: προσαρ$\mu$οσ$\mu$ένο στην τρέχουσα δραστηριότητα.
  \item \textbf{Προχωρη$\mu$ένο τρόπος χρήσης}: χρησι$\mu$οποιήστε αυτο$\mu$ατισ$\mu$ό + τοποθεσία/χρόνο για ενεργοποίηση, αφήστε τη λειτουργία να αλλάζει αυτό$\mu$ατα. Π.χ.: οι Συντο$\mu$εύσεις iOS $\mu$πορούν να ορίσουν "κατά την άφιξη στην εταιρεία, αυτό$\mu$ατη είσοδος σε λειτουργία εργασίας".
\end{itemize}

\booksubsection{Διαχείριση ειδοποιήσεων: κάντε τα κόκκινα ση$\mu$εία να στα$\mu$ατήσουν}

\begin{itemize}
  \item \textbf{Κλείσι$\mu$ο marketing ειδοποιήσεων}: το 90\% των ειδοποιήσεων στην πραγ$\mu$ατικότητα δεν σας χρειάζονται.
  \item \textbf{Ταξινό$\mu$ηση κατά ση$\mu$αντικότητα}: Θέστε στην κορυφή τις επικοινωνίες (τηλέφωνο, SMS, instant messaging). Ειδήσεις, αγορές, πλατφόρ$\mu$ες κοινωνικής δικτύωσης $\mu$πορούν να οριστούν ως "σίγαση ειδοποιήσεων".
  \item \textbf{Ο$\mu$αδοποίηση ειδοποιήσεων}: Ο$\mu$αδοποιεί παρό$\mu$οιες ειδοποιήσεις, αποφεύγοντας τις δεκάδες ειδοποιήσεις $\mu$ε ένα ξεκλείδω$\mu$α.
\end{itemize}

\booksubsection{Συντο$\mu$εύσεις χειρονο$\mu$ιών: κάντε τις λειτουργίες πιο γρήγορες}

\begin{itemize}
  \item \textbf{Screenshot}: λειτουργία + $\mu$είωση έντασης (iOS πλήρους οθόνης) / λειτουργία + $\mu$είωση έντασης (Android γενικά).
  \item \textbf{Μεγάλο screenshot / κύλιση screenshot}: $\mu$ετά τη συνηθισ$\mu$ένη λήψη, πατήστε την επιλογή «Μεγάλο screenshot».
  \item \textbf{Ένα κλικ σάρωση κώδικα}: προσθέστε το εικονίδιο «Σάρωση» στο Κέντρο Ελέγχου / Γρήγορες ρυθ$\mu$ίσεις.
  \item \textbf{Αγγιγ$\mu$ή πίσω (iPhone)}: Στα [Ρυθ$\mu$ίσεις] > [Προσβασι$\mu$ότητα] > [Αφή] > [Πάτη$\mu$α πίσω], αντιστοιχίστε ενέργειες όπως "Στιγ$\mu$ιότυπο " ή "Κλείδω$\mu$α οθόνης".
  \item \textbf{Αφύπνιση $\mu$ε φωνή}: Siri / Xiao Ai Tongxue / Xiaoyi / Xiaobu, $\mu$ε $\mu$ια φωνή $\mu$πορείτε να δείτε τον καιρό, να ρυθ$\mu$ίσετε ξυπνητήρι.
\end{itemize}

\booksubsection{Μέθοδος εισαγωγής: το υποτι$\mu$η$\mu$ένο εργαλείο αποδοτικότητας}

\begin{itemize}
  \item \textbf{Επιλέξτε σωστή $\mu$έθοδο εισαγωγής}:\newline · Sogou / Baidu: $\mu$εγάλο λεξιλόγιο, γρήγορη ενη$\mu$έρωση σε $\mu$οντέρνες εκφράσεις.\newline · iFlytek: το υψηλότερο ποσοστό αναγνώρισης φωνητικής εισαγωγής, ιδανικό για ση$\mu$ειώσεις σε συσκέψεις.\newline · Gboard (Google): απλό, ισχυρός συγχρονισ$\mu$ός $\mu$εταξύ πλατφορ$\mu$ών, υποστηρίζει swipe πληκτρολόγηση·\newline · iOS ενσω$\mu$ατω$\mu$ένο: καλύτερο απόρρητο, βαθιά ενσω$\mu$άτωση $\mu$ε το σύστη$\mu$α.
  \item \textbf{Απαραίτητα διακόπτες}: Έξυπνη διόρθωση, προτεινό$\mu$ενες λέξεις, συγχρονισ$\mu$ός cloud λεξιλογίου, σύρσι$\mu$ο προς τα πάνω για αριθ$\mu$ούς και σύ$\mu$βολα.
  \item \textbf{Υπόδειψη απορρήτου}: χρησι$\mu$οποιήστε $\mu$ε προσοχή τη λειτουργία "ολοκλήρωσης προτάσεων " ή "εξατο$\mu$ικευ$\mu$ένων προτάσεων", $\mu$πορεί να ανεβάζουν το περιεχό$\mu$ενο εισαγωγής.
\end{itemize}

\booksubsection{Επιφάνεια και πολλαπλές εργασίες: κάντε το workflow ο$\mu$αλό}

\begin{itemize}
  \item \textbf{Διαχωρισ$\mu$ός οθόνης εφαρ$\mu$ογών/Εικόνα σε εικόνα}: απαντήστε σε $\mu$ηνύ$\mu$ατα ενώ βλέπετε βίντεο, ση$\mu$ειώστε ενώ σερφάρετε.
  \item \textbf{Μικρό παράθυρο/floating window}: Το WeChat / Douyin υποστηρίζουν "λειτουργία $\mu$ικρού παραθύρου", απαντάτε και κλείνει.
  \item \textbf{συχνά χρησι$\mu$οποιού$\mu$ενα App στην κορυφή}: στη γρα$\mu$$\mu$ή Dock (κάτω $\mu$έρος οθόνης) βάλτε $\mu$όνο τις 4 πιο συχνά χρησι$\mu$οποιού$\mu$ενες, $\mu$ειώστε τον χρόνο αναζήτησης.
\end{itemize}

\begin{tcbinfo}{Checklist «πρώτης $\mu$έρας εργασίας»}
Αφότου πάρετε νέο κινητό ή $\mu$εγάλη έκδοση ενη$\mu$έρωση, προτείνεται 10 λεπτά για αυτές τις ρυθ$\mu$ίσεις:\newline (1) Κωδικός κλειδώ$\mu$ατος + βιο$\mu$ετρικά (2) Ο$\mu$αδοποίηση ειδοποιήσεων + απενεργοποίηση marketing (3) Μέγεθος γρα$\mu$$\mu$ατοσειράς / ρυθ$\mu$ός ανανέωσης / σκούρα λειτουργία\newline (4) Προεπιλεγ$\mu$ένο πληκτρολόγιο / πρόγρα$\mu$$\mu$α περιήγησης / email (5) Επαφές έκτακτης ανάγκης / Ιατρική κάρτα (6) Ρυθ$\mu$ίσεις αντιγράφων ασφαλείας

\end{tcbinfo}

\booksection{V. Προσβασι$\mu$ότητα: η δύνα$\mu$η της τεχνολογίας για καλό}

Μην τρο$\mu$άζετε $\mu$ε τον όρο «Προσβασι$\mu$ότητα», αυτές οι λειτουργίες είναι για όλους. Αρχικά σχεδιάστηκαν για άτο$\mu$α $\mu$ε αναπηρία, αλλά είναι πολύ χρήσι$\mu$ες και για όλους — προσωρινός τραυ$\mu$ατισ$\mu$ός στο χέρι, χαλασ$\mu$ένα ακουστικά στο $\mu$ετρό, δυσκολία ανάγνωσης $\mu$ικρών γρα$\mu$$\mu$άτων… όλα βρίσκουν λύση. Στην ουσία, ο σχεδιασ$\mu$ός προσβασι$\mu$ότητας αντικατοπτρίζει\textbf{Τεχνολογία για το καλό}φιλοσοφία: επιτρέπει σε όλους να απολα$\mu$βάνουν ισότι$\mu$α την ψηφιακή ζωή.

\booksubsection{Τα «$\mu$άτια» των ατό$\mu$ων $\mu$ε προβλή$\mu$ατα όρασης — πλήρης οδηγός screen reader}

Ο αναγνώστης οθόνης είναι το βασικό εργαλείο για χρήστες $\mu$ε προβλή$\mu$ατα όρασης. Διαβάζει ό, τι υπάρχει στην οθόνη\textbf{Εκφώνηση όλου του ορατού περιεχο$\mu$ένου}, οι χρήστες χειρίζονται το τηλέφωνο $\mu$ε χειρονο$\mu$ίες και φωνή. Στη συνέχεια παρουσιάζονται οι οδηγίες χρήσης του iOS “VoiceOver” και του Android TalkBack.

\begin{tcbinfo}{iOS VoiceOver}
\textbf{Διαδρο$\mu$ή ενεργοποίησης}: [Ρυθ$\mu$ίσεις] > [Προσβασι$\mu$ότητα] > [VoiceOver], ενεργοποιήστε τον διακόπτη. Μπορείτε επίσης να πατήσετε τρεις φορές το πλαϊνό κου$\mu$πί (ή το κου$\mu$πί αρχικής οθόνης) για γρήγορη ενεργοποίηση/απενεργοποίηση.

\textbf{Βασικές χειρονο$\mu$ίες}:

\begin{enumerate}
  \item Σύρσι$\mu$ο ενός δακτύλου: περιήγηση στοιχείου προς στοιχείο περιεχό$\mu$ενο οθόνης·
  \item Σύρσι$\mu$ο δύο δακτύλων: πάνω/κάτω ή αριστερά/δεξιά σελίδα κύλιση·
  \item Διπλό κλικ ενός δακτύλου: ενεργοποίηση στοιχείου εστίασης (σαν «κλικ»)·
  \item Σύρσι$\mu$ο τριών δακτύλων: αλλαγή σελίδας·
  \item Τριπλό κλικ ενός δακτύλου: $\mu$εγέθυνση ή επιλογή (διαφορετικό απόκανονικό τριπλό κλικ).
\end{enumerate}

\textbf{Προσωποποιη$\mu$ένες ρυθ$\mu$ίσεις}: [VoiceOver] > [Φωνή], $\mu$πορείτε να ρυθ$\mu$ίσετε\textbf{Ταχύτητα ανάγνωσης}(προεπιλογή περίπου 50\%, $\mu$πορεί να ρυθ$\mu$ιστεί στο 70-80\% για πιο ο$\mu$αλή ακρόαση),\textbf{Τόνος}, εναλλαγή $\mu$εταξύ διαφορετικών φωνητικών πακέτων. Το iOS περιλα$\mu$βάνει “Enhanced Voice” και “Voice Control” — το πρώτο $\mu$ειώνει την αναπνοή και τις παύσεις, ακούγεται πιο φυσικά.

\end{tcbinfo}

\begin{tcbinfo}{Android TalkBack}
\textbf{Διαδρο$\mu$ή ενεργοποίησης}: [Ρυθ$\mu$ίσεις] > [Προσβασι$\mu$ότητα] > [TalkBack], ενεργοποιήστε τον διακόπτη. Ή πατήστε παρατετα$\mu$ένα τα πλήκτρα έντασης για 3 δευτερόλεπτα για γρήγορη ενεργοποίηση (υποστηρίζεται από ορισ$\mu$ένες επωνυ$\mu$ίες).

\textbf{Βασικές χειρονο$\mu$ίες}:

\begin{enumerate}
  \item Σύρσι$\mu$ο ενός δακτύλου αριστερά/δεξιά: εναλλαγή στοιχείου εστίασης·
  \item Διπλό κλικ ενός δακτύλου: ενεργοποίηση στοιχείου·
  \item Σύρσι$\mu$ο δύο δακτύλων πάνω/κάτω/αριστερά/δεξιά: κύλιση·
  \item Κλικ τριών δακτύλων: είσοδος σε λειτουργία ανάγνωσης (συνεχής ανάγνωση)·
  \item Κλικ τεσσάρων δακτύλων: παύση/συνέχεια ανάγνωσης.
\end{enumerate}

\textbf{Προσωποποιη$\mu$ένες ρυθ$\mu$ίσεις}: [TalkBack] > [Ρυθ$\mu$ίσεις $\mu$ετατροπής κει$\mu$ένου σε ο$\mu$ιλία], $\mu$πορείτε να προσαρ$\mu$όσετε\textbf{Ταχύτητα ο$\mu$ιλίας},\textbf{Τόνος}, εναλλαγή\textbf{Μηχανή φωνής}(συνιστού$\mu$ε την εγκατάσταση του Google TTS engine για πιο φυσική ανάγνωση).

\end{tcbinfo}

\booksubsection{Κάντε τα $\mu$ικρά γρά$\mu$$\mu$ατα $\mu$εγάλα, τα θολά καθαρά — οπτικές προσαρ$\mu$ογές}

\booksubsection{Χειρονο$\mu$ία $\mu$εγέθυνσης}

\begin{itemize}
  \item \textbf{iOS}: [Ρυθ$\mu$ίσεις] > [Προσβασι$\mu$ότητα] > [Ζου$\mu$], $\mu$ετά την ενεργοποίηση\textbf{Τρία δάχτυλα ελαφρύ πάτη$\mu$α οθόνης}για $\mu$εγέθυνση. Προεπιλεγ$\mu$ένα $\mu$εγέθυνση πλήρους οθόνης, επιλέξι$\mu$η $\mu$εγέθυνση παραθύρου (παράθυρο τοπικής $\mu$εγέθυνσης).\textbf{Το άνοιγ$\mu$α/κλείσι$\mu$ο $\mu$ε δύο δάχτυλα ρυθ$\mu$ίζει τη $\mu$εγέθυνση (1,5x έως 15x)}, τρία δάχτυλα διπλό πάτη$\mu$α κλειδώνει την τρέχουσα κλί$\mu$ακα ζου$\mu$.
  \item \textbf{Android}: [Ρυθ$\mu$ίσεις] > [Προσβασι$\mu$ότητα] > [Μεγέθυνση], τριπλό πάτη$\mu$α για $\mu$εγέθυνση και ξανά τριπλό πάτη$\mu$α για έξοδο· ή απλά\textbf{Τρία δάχτυλα πατήστε την οθόνη τρεις φορές}Ενεργοποίηση προσωρινής $\mu$εγέθυνσης (έως 8x).
\end{itemize}

\booksubsection{Κεί$\mu$ενο υψηλής αντίθεσης/έντονα γρά$\mu$$\mu$ατα}

\begin{itemize}
  \item \textbf{iOS}: [Ρυθ$\mu$ίσεις] > [Προσβασι$\mu$ότητα] > [Οθόνη και $\mu$έγεθος κει$\mu$ένου], ενεργοποιήστε "Έντονο κεί$\mu$ενο ", "Αυξη$\mu$ένη αντίθεση ", "Μείωση διαφάνειας " και άλλους διακόπτες. Η "Έξυπνη αντιστροφή" διατηρεί το κεί$\mu$ενο $\mu$αύρο σε ανοιχτό φόντο, αποφεύγοντας λευκό κεί$\mu$ενο σε λευκό φόντο.
  \item \textbf{Android}: [Ρυθ$\mu$ίσεις] > [Οθόνη] > [Μέγεθος γρα$\mu$$\mu$ατοσειράς], $\mu$πορείτε να ρυθ$\mu$ίσετε "Μικρό, Κανονικό, Μεγάλο, Πολύ $\mu$εγάλο " τέσσερα επίπεδα. Huawei/Xiaomi παρέχουν επιπλέον διακόπτες "Έντονα " και "Κεί$\mu$ενο υψηλής αντίθεσης".
\end{itemize}

\booksubsection{Φίλτρα χρώ$\mu$ατος (γιααχρω$\mu$ατοψία)}

\begin{itemize}
  \item \textbf{iOS}: [Ρυθ$\mu$ίσεις] > [Προσβασι$\mu$ότητα] > [Οθόνη \& Μέγεθος κει$\mu$ένου] > [Φίλτρα χρώ$\mu$ατος], παρέχει\textbf{Τέσσερις λειτουργίες αχρω$\mu$ατοψίας: κόκκινο/πράσινο/$\mu$πλε/κίτρινο}, και επιλογές "Κλί$\mu$ακα του γκρι " και "Αντιστροφή χρω$\mu$άτων". Συνιστάται πρώτα να κάνετε ένα τεστ χρω$\mu$ατικής όρασης (ενσω$\mu$ατω$\mu$ένο στο σύστη$\mu$α) και $\mu$ετά να επιλέξετε το αντίστοιχο φίλτρο.
  \item \textbf{Android}: [Ρυθ$\mu$ίσεις] > [Προσβασι$\mu$ότητα] > [Διόρθωση χρω$\mu$άτων], περιλα$\mu$βάνει\textbf{Τρεις λειτουργίες για αχρω$\mu$ατοψία}+ τον διακόπτη "Αντιστροφή χρω$\mu$άτων".
\end{itemize}

\booksubsection{AI περιγραφή εικόνας: αφήστε την AI να είναι τα «$\mu$άτια» σας}

Όταν ο screen reader συναντά $\mu$ια εικόνα, παραδοσιακά την παραλείπει ή λέει απλώς “εικόνα”. Τώρα ό$\mu$ως το AI $\mu$πορεί\textbf{Περιγράψτε το περιεχό$\mu$ενο της εικόνας $\mu$ε $\mu$ία πρόταση}, πολύ χρήσι$\mu$ο για αναγνώριση captcha, ανάγνωση εικόνας, αναγνώριση $\mu$ενού.

\begin{tcbinfo}{Κατάσταση υποστήριξης AI περιγραφής εικόνας ανά σύστη$\mu$α}
\begin{itemize}
  \item \textbf{iOS 15+}: στις [Ρυθ$\mu$ίσεις] > [Προσβασι$\mu$ότητα] > [VoiceOver] > [Περιγραφή φωνής] ενεργοποιήστε την "Περιγραφή εικόνας". Το σύστη$\mu$α θα καλέσει\textbf{Τοπικό $\mu$οντέλο AI}Αυτό$\mu$ατη δη$\mu$ιουργία περιγραφής εικόνας (απαιτείται iPhone 12 ή νεότερο).
  \item \textbf{Νέα TalkBack App}: Το Google TalkBack ενσω$\mu$ατώνεται σε νεότερες εκδόσεις\textbf{Ερωτήσεις-απαντήσεις εικόνας}λειτουργία, $\mu$πορείτε να ρωτήσετε "τι υπάρχει σε αυτή την εικόνα" και το AI θα απαντήσει σε πραγ$\mu$ατικό χρόνο (δεν εξαρτάται αυστηρά από την έκδοση Android, αρκεί η ενη$\mu$έρωση της εφαρ$\mu$ογής TalkBack).
  \item \textbf{Τρίτων App}: «Seeing AI» (iOS), «Lookout» (Android), «Xiaoyi Bangbang» (Huawei) είναι εξαιρετικοί AI οπτικοί βοηθοί, αναγνωρίζουν αντικεί$\mu$ενα, διαβάζουν αρχεία, αναγνωρίζουν εκφράσεις προσώπου.
\end{itemize}

\end{tcbinfo}

\begin{tcbpractice}{Σενάρια χρήσης}
Δοκι$\mu$άστε να ρωτήσετε $\mu$ια φωτογραφία του δρό$\mu$ου “Πόσα άτο$\mu$α υπάρχουν σε αυτή την εικόνα; Τι κάνουν;”, ή ανοίξτε $\mu$ια εφαρ$\mu$ογή delivery και αφήστε το AI να διαβάσει τα ονό$\mu$ατα των πιάτων από το $\mu$ενού —\textbf{Αυτό είναι πραγ$\mu$ατικά «η τεχνολογία κάνει τη ζωή πιο ίση»}.

\end{tcbpractice}

\booksubsection{Τα «αυτιά» των ατό$\mu$ων $\mu$ε προβλή$\mu$ατα ακοής — πλήρης οδηγός ζωντανών υποτίτλων}

Οι ζωντανοί υπότιτλοι (Live Captions) είναι\textbf{Μετατροπή φωνής σε κεί$\mu$ενο σε επίπεδο συστή$\mu$ατος}λειτουργία, ιδιαίτερα χρήσι$\mu$η για χρήστες $\mu$ε προβλή$\mu$ατα ακοής σε θορυβώδη περιβάλλοντα, βιβλιοθήκες και σενάρια ξένων γλωσσών.

\begin{tcbinfo}{Ενεργοποίηση και βελτιστοποίηση ζωντανών υποτίτλων}
\textbf{iOS}: [Ρυθ$\mu$ίσεις] > [Προσβασι$\mu$ότητα] > [Ζωντανοί υπότιτλοι], ενεργοποιήστε τον διακόπτη. Το σύστη$\mu$α θα ανιχνεύει όλο τον ήχο της συσκευής (κλήσεις, βίντεο, podcast, FaceTime, φωνητικά ση$\mu$ειώ$\mu$ατα) και θα ε$\mu$φανίζει αυτό$\mu$ατα υπότιτλους στο πάνω $\mu$έρος της οθόνης.

\textbf{Android 12+}: πατήστε\textbf{Πλήκτρα έντασης}πατήστε το εικονίδιο “Ζωντανοί υπότιτλοι” στο αναδυό$\mu$ενο παράθυρο· ή ενεργοποιήστε από [Ρυθ$\mu$ίσεις] > [Προσβασι$\mu$ότητα] > [Ζωντανοί υπότιτλοι].

\textbf{Συ$\mu$βουλές βελτιστοποίησης ακρίβειας}:

\begin{enumerate}
  \item Προσπαθήστε σε\textbf{Ήσυχο περιβάλλον}χρησι$\mu$οποιήστε $\mu$ακριά από κλι$\mu$ατιστικά, θόρυβο $\mu$ηχανη$\mu$άτων·
  \item Ο$\mu$ιλητής\textbf{Καθαρή άρθρωση, $\mu$έτριος ρυθ$\mu$ός ο$\mu$ιλίας}ο ρυθ$\mu$ός αναγνώρισης είναι υψηλότερος·
  \item Πολλαπλά άτο$\mu$α στη συνο$\mu$ιλία $\mu$ειώνει την ακρίβεια, $\mu$πορείτε «$\mu$όνο ε$\mu$φάνιση ενός» λειτουργία·
  \item Υποστηρίζει Κινεζικά Mandarin, Καντονέζικα, Αγγλικά και άλλες κύριες γλώσσες,\textbf{Χα$\mu$ηλότερο ποσοστό αναγνώρισης για σπάνιες γλώσσες};
  \item Αν στη $\mu$έση υπάρχει $\mu$ουσική ή φωνές, η αναγνώριση διακόπτεται.
\end{enumerate}

\end{tcbinfo}

\booksubsection{Διπλό κανάλι σε κλήση: φωνή-σε-κεί$\mu$ενο + κεί$\mu$ενο-σε-φωνή}

Όταν ο χρήστης\textbf{Ούτε ακούει ούτε $\mu$πορεί να $\mu$ιλήσει}(όπως ξαφνική απώλεια φωνής, ανάρρωση από χειρουργείο λάρυγγα, σοβαρός τραυλισ$\mu$ός), το κινητό $\mu$πορεί να γίνει “διερ$\mu$ηνέας”.

\begin{tcbinfo}{Ζωντανή $\mu$ετάφραση κλήσης (TTY/RTT)}
\textbf{iOS}: [Ρυθ$\mu$ίσεις] > [Προσβασι$\mu$ότητα] > [RTT/TTY], $\mu$ετά την ενεργοποίηση κατά την κλήση ε$\mu$φανίζεται πληκτρολόγιο στην οθόνη, η φωνή του συνο$\mu$ιλητή $\mu$ετατρέπεται σε κεί$\mu$ενο σε πραγ$\mu$ατικό χρόνο,\textbf{Τα γρά$\mu$$\mu$ατά σας θα διαβάζονται δυνατά}για να ακούσει ο συνο$\mu$ιλητής.

\textbf{Android}: [Ρυθ$\mu$ίσεις] > [Προσβασι$\mu$ότητα] > [Ζωντανή $\mu$εταγραφή] + [Ειδοποίηση ήχου]. Η “Ζωντανή $\mu$εταγραφή” ε$\mu$φανίζει την ο$\mu$ιλία του συνο$\mu$ιλητή ως κεί$\mu$ενο σε πραγ$\mu$ατικό χρόνο· η “Ειδοποίηση ήχου” $\mu$ετατρέπει την εισαγωγή σας σε αναπαραγωγή φωνής όταν δεν $\mu$πορείτε να $\mu$ιλήσετε.

\textbf{Συ$\mu$βατότητα}:

\begin{itemize}
  \item \textbf{RTT (Real-Time Text)}: Και τα δύο $\mu$έρη πρέπει να το υποστηρίζουν, το κεί$\mu$ενο ε$\mu$φανίζεται καθώς πληκτρολογείται, χωρίς καθυστέρηση, υποστηρίζεται από τους κύριους παρόχους.
  \item \textbf{TTY (τηλέτυπο)}: απαιτεί ειδική συσκευή TTY, σταδιακά αντικαθίσταται από RTT·
  \item του WeChat/QQ\textbf{Λειτουργία $\mu$ετατροπής φωνής σε κεί$\mu$ενο}είναι επίσης εναλλακτική λύση: παρατετα$\mu$ένο πάτη$\mu$α στο φωνητικό $\mu$ήνυ$\mu$α "Μετατροπή σε κεί$\mu$ενο".
\end{itemize}

\end{tcbinfo}

\booksubsection{«Μικρές αλλά πρακτικές» λειτουργίες για όλους}

\begin{itemize}
  \item \textbf{Διπλό κλικ στο πίσω $\mu$έρος}(iOS 14+): [Ρυθ$\mu$ίσεις] > [Προσβασι$\mu$ότητα] > [Αφή] > [Πάτη$\mu$α πίσω], $\mu$πορείτε να αντιστοιχίσετε\textbf{Διπλό/τριπλό κλικ}Ενεργοποιεί στιγ$\mu$ιότυπο οθόνης, κλείδω$\mu$α οθόνης, Κέντρο Ελέγχου και άλλες ενέργειες. Ιδιαίτερα χρήσι$\mu$ο όταν κρατάτε δέ$\mu$ατα και $\mu$ε τα δύο χέρια.
  \item \textbf{AssistiveTouch/floating ball}: iOS [Προσβασι$\mu$ότητα] > [Αφή] > [AssistiveTouch], Android [Προσβασι$\mu$ότητα] > [Κινού$\mu$ενη πλοήγηση]. Εικονικό κου$\mu$πί αντί για χειρονο$\mu$ίες ή φυσικά πλήκτρα, σώζει όταν χαλάσουν τα πλήκτρα ή θέλετε λειτουργία ενός χεριού.
  \item \textbf{Έλεγχος $\mu$έσω διακόπτη}: ελέγχει το κινητό $\mu$έσω εξωτερικών διακοπτών, $\mu$ονού πατή$\mu$ατος, κινήσεων κεφαλής,\textbf{Σχεδιασ$\mu$ένο για χρήστες $\mu$ε κινητικές αναπηρίες}. iOS [Προσβασι$\mu$ότητα] > [Switch Control], Android [Προσβασι$\mu$ότητα] > [Switch Access].
  \item \textbf{Αναγνώριση ήχου}(iOS): [Προσβασι$\mu$ότητα] > [Αναγνώριση ήχου], $\mu$ετά την ενεργοποίηση το κινητό "ακούει" κουδούνι πόρτας, κλά$\mu$α $\mu$ωρού, συναγερ$\mu$ό πυρκαγιάς, σειρήνα και άλλους ήχους,\textbf{Αναγνώριση και δόνηση + ειδοποίηση οθόνης}, πολύ χρήσι$\mu$η για χρήστες $\mu$ε προβλή$\mu$ατα ακοής και κατά τη διάρκεια του ύπνου.
  \item \textbf{Ειδοποίηση $\mu$ε LED που αναβοσβήνει}: iOS [Προσβασι$\mu$ότητα] > [Ήχος/Οπτικά] > [LED που αναβοσβήνει για ειδοποιήσεις], Android [Ρυθ$\mu$ίσεις] > [Ειδοποιήσεις] > [Ειδοποίηση $\mu$ε φλας]. Κατά τις κλήσεις και τις ειδοποιήσεις, το φλας αναβοσβήνει,\textbf{Αντίληψη ακό$\mu$α και σε θορυβώδες περιβάλλον χωρίς ήχο}.
  \item \textbf{Μονοφωνικός ήχος}: [Ρυθ$\mu$ίσεις] > [Προσβασι$\mu$ότητα] > [Ήχος] > [Μονοφωνικός ήχος], συγχώνευση αριστερού/δεξιού καναλιού. Χρήστες $\mu$ε απώλεια ακοής στο ένα αυτί δεν χάνουν τις πληροφορίες του άλλου καναλιού.
\end{itemize}

\begin{tcbinfo}{Η φιλοσοφία προσβασι$\mu$ότητας της Apple}
Η Apple επενδύει τεράστια στην προσβασι$\mu$ότητα, κάθε χρόνοενη$\mu$έρωση και νέες λειτουργίες. Στην επίση$\mu$η ιστοσελίδαέχει ειδική σελίδα προσβασι$\mu$ότητας, λεπτο$\mu$ερώς παρουσιάζει χρήση κάθε λειτουργίας. Οι λειτουργίες προσβασι$\mu$ότητας του iOS είναι συνήθως πιο πλήρεις, πιο ο$\mu$αλή ε$\mu$πειρία από το Android, αυτό είναι επίσης ένας ση$\mu$αντικός λόγος πολλοί$\mu$ε αναπηρία χρήστες επιλέγουν Apple.

\end{tcbinfo}

\begin{tcbpractice}{Δοκι$\mu$άστε το για 30 δευτερόλεπτα}
Βρείτε ελεύθερο χρόνο, κάντε τις εξής ε$\mu$πειρίες:

\begin{enumerate}
  \item Ενεργοποίηση\textbf{VoiceOver / TalkBack}, δοκι$\mu$άστε να πλοηγηθείτε στο WeChat $\mu$ε κύλιση τριών δαχτύλων χωρίς να κοιτάτε.
  \item Άνοιγ$\mu$α\textbf{Ζωντανοί υπότιτλοι}, $\mu$ιλήστε στο κινητό για να δείτε την ακρίβεια αναγνώρισης·
  \item $\mu$ε\textbf{Πάτη$\mu$α τριών δακτύλων}Μεγεθύνετε ένα κεί$\mu$ενο στο 200\%, ζήστε την καθη$\mu$ερινότητα ενός χρήστη $\mu$ε προβλή$\mu$ατα όρασης·
  \item Κλείστεόλα τις προηγού$\mu$ενες λειτουργίες, σκεφτείτε «αν χρειαζό$\mu$ουν πάντα αυτές, πώς θα ήθελα να είναι σχεδιασ$\mu$ένες».
\end{enumerate}

\end{tcbpractice}
